\chapter{Building an incremental request parser}
\chapterintro{Parsing turns hostile bytes into a trusted internal structure. The parser's most important output is not convenience; it is one unambiguous interpretation of exactly one message.}

\begin{objectives}
\begin{itemize}
  \item Separate incremental buffering, syntactic parsing, framing, and semantic validation.
  \item Build structured requests without losing repeated fields or trailing bytes.
  \item Reject malformed lines, unsupported versions, conflicting framing, and invalid targets.
  \item Test the parser independently of sockets with fragmented inputs.
\end{itemize}
\end{objectives}

\section{From a teaching parser to a strict parser}

Stage 3 introduces a readable request dataclass and a direct parsing function.
It is intentionally smaller than the final implementation.

\codefile{../code/stages/03_request_parser/server.py}

The stage establishes the right order: find the header boundary, parse the start
line and fields, determine content length, receive the declared body, then create
a request. It still uses a one-value dictionary for fields, does not preserve
leftover bytes, and handles errors with broad 400 responses. Those limitations
motivate the final types rather than remaining hidden.

\section{The final parsing contract}

\code{RequestParser} owns one byte buffer for one TCP connection. Calling
\code{feed} never promises a complete request. Calling \code{pop\_request}
returns either one structured request or \code{None}. A protocol violation raises
\code{HTTPError} with a safe status and detail.

The algorithm is:

\begin{enumerate}
  \item find the first CRLF--CRLF header boundary;
  \item if absent, apply early line/header limits and request more bytes;
  \item validate strict line endings;
  \item parse the request line and require HTTP/1.1;
  \item parse each field into a multi-value collection and require one Host;
  \item determine body length or reject unsupported/ambiguous framing;
  \item if the declared body is incomplete, return \code{None};
  \item copy exactly the body, delete exactly the consumed prefix, parse the
        target, and construct \code{Request}.
\end{enumerate}

This order avoids interpreting a partial body and preserves a following request.

\section{Read the complete final parser}

\codefile{../code/final/tinyhttp/parser.py}

\subsection{Tokens and line syntax}

The method and field name use a token allow-list. The request line must contain
exactly three non-empty space-separated components. This intentionally rejects
tabs and extra spaces instead of trying to guess intent. Version mismatches
receive 505 rather than falling through into a parser designed for different
syntax.

Header values are decoded with ISO-8859-1 as a byte-preserving mapping for the
HTTP/1.x field octets accepted by this implementation. Individual application
fields still need their own semantics. Control characters and obsolete folding
are rejected.

\subsection{Framing before content}

The parser inspects every occurrence of \code{Transfer-Encoding} and
\code{Content-Length}. It never accepts both. It rejects repeated lengths and
does not wait for a body whose declared size already exceeds the limit. Only a
single valid decimal length is used to slice the stream.

This decision belongs in the HTTP parser, not in a JSON handler. A handler should
receive only a body whose byte boundary is already established.

\subsection{Targets and decoding}

The parser accepts a path beginning with slash or the asterisk target. It rejects
an absolute URI because this server is an origin, not a forward proxy. It rejects
fragments, which are client-side URI components and should not be sent as HTTP
request targets.

Percent escapes are validated before decoding. Path segments are decoded
separately so an encoded slash inside a segment does not secretly create a new
routing segment. Query parsing preserves repeated values in lists. Invalid UTF-8
becomes a 400 problem instead of leaking a Unicode exception.

\section{Internal trust is earned in layers}

After parser success, code may trust that the request has one message boundary,
a legal supported version, valid field syntax, and a parsed target. It may not
yet trust that JSON is well formed, a task status is allowed, a referenced task
exists, or a user is authorized. Each layer narrows uncertainty according to
its expertise.

\begin{mentalmodel}
Validation is not one function at the edge. Transport validates stream progress;
HTTP validates message syntax and framing; content code validates media syntax;
the application validates domain rules; storage enforces durable constraints.
\end{mentalmodel}

\section{Parser tests as executable grammar examples}

The test suite feeds requests one byte at a time, puts two requests in one feed,
tests repeated fields, rejects missing Host, rejects transfer coding, and applies
limits before waiting for an oversized body. These cases are more valuable than
testing only a browser's ordinary GET.

For every accepted example, add fragmentation variants. For every rejected
example, assert both status and reason category. A parser test should not depend
on timing or the operating system's packet choices.

\begin{tryit}
From \filepath{code/final}, run:
\begin{wiremessage}
python -m unittest tests.test_models_parser -v
\end{wiremessage}
Then add a test for a malformed percent escape such as \code{/tasks/\%ZZ}. Make
the expectation explicit before running it.
\end{tryit}

\summarytitle

The incremental parser owns a connection buffer and consumes exactly one framed
message. It validates strict start/field syntax, preserves repeated fields,
rejects ambiguous framing, enforces limits early, and produces decoded path and
query structures. Later layers receive a safer request, not magically trusted
application data.

\reviewtitle
\begin{enumerate}
  \item Why can \code{pop\_request} return \code{None} without indicating an error?
  \item Which bytes remain after one request is removed from the buffer?
  \item Why are percent escapes checked before decoding?
  \item What validation still remains after HTTP parsing succeeds?
\end{enumerate}

\begin{exercise}[Adversarial table]
Create a table of at least twelve requests: four valid, four incomplete, and four
invalid. Include duplicate lengths, invalid Host, a large declared body, an
encoded path, and two buffered requests. Record expected parser outcome and
status where applicable.
\end{exercise}

\begin{exercise}[Invariant proof]
State an invariant for the parser buffer before and after \code{pop\_request}.
Explain why deleting exactly \code{header\_end + content\_length} bytes preserves
that invariant for persistent connections.
\end{exercise}
